\documentclass{situationist-preprint}

\title{Thought Experiments as Arguments}
\msid{SCI-15-a}
\msfamily{SCI-15}
\msdiscipline{Philosophy of Science}
\mstype{Principal manuscript}
\mskeywords{thought experiments; Norton; Brown; Galileo; falling bodies; empiricism; mental models; argument}

\begin{document}
\maketitle

\begin{abstract}
Thought experiments seem to teach us about nature without new observation. Galileo's argument that heavy and light bodies fall at the same rate is the standard example. \citet{Brown1991} holds that some thought experiments give a priori insight into the laws of nature; \citet{Norton1996,Norton2004} holds that they are arguments in picturesque form. We defend Norton's view in a specific form: thought experiments can teach us about nature, and their evidential force is exactly that of the arguments they encode. The key premise is that a thought experiment supplies no new information about the world, so that whatever justification its conclusion has derives from information the thinker already possesses, transformed by a process whose reliability is the reliability of an inference. We show that the evidential force of a thought experiment neither exceeds nor falls short of that of its reconstruction, apply the account to Galileo's falling bodies, and answer the objections from platonism, from tacit knowledge and from mental models.
\end{abstract}

\section{The problem}

A thought experiment describes an imagined scenario and invites the reader to judge what would happen in it, and the judgement is meant to support a conclusion about the actual world. In the \emph{Discorsi}, Galileo asks us to suppose, with Aristotle, that heavier bodies fall faster, and to imagine a heavy body tied to a light one (Galileo, \emph{Two New Sciences}, First Day). The light body should retard the heavy one, so the pair should fall more slowly than the heavy body alone; but the pair is heavier than the heavy body, so it should fall faster. The supposition is contradictory, and Galileo concludes that bodies of different weight fall at the same rate. No one dropped anything, yet a conclusion about falling bodies seems to have been established.

How can this be? \citet{Brown1991} answers that some thought experiments, which he calls platonic, give us a priori access to laws of nature, understood as relations among abstract universals that the mind can grasp. \citet{Norton1996} answers that thought experiments are arguments: they posit hypothetical or counterfactual states of affairs and invoke particulars irrelevant to the generality of the conclusion, but their epistemic work is done by an argument from premises the thinker already accepts. \citet{Gendler1998} holds an intermediate view, on which some thought experiments justify by bringing to bear tacit knowledge that the explicit argument does not capture, and \citet{Nersessian1993} and \citet{Miscevic1992} hold that thought experiments are manipulations of mental models rather than arguments. \citet{Kuhn1977} held that thought experiments can teach by revealing conceptual tensions that force a reconceptualization.

We argue that thought experiments can teach us about nature, and that their evidential force is that of the arguments they encode.

\section{Framework}

\begin{definition}[Thought experiment]\label{d:te}
A \emph{thought experiment} is a process in which a thinker considers an imagined scenario, judges what would occur in it, and draws from that judgement a conclusion about the actual world.
\end{definition}

\begin{definition}[Encoded argument]\label{d:enc}
An argument \emph{encodes} a thought experiment if its premises are propositions that the thinker accepts and on which her judgement about the scenario depends, its conclusion is the conclusion of the thought experiment, and its inferential steps correspond to the steps by which the judgement is reached.
\end{definition}

The premises of an encoded argument may be tacit, in the sense that the thinker does not formulate them, provided her judgement depends on them in the sense that she would revise it if she ceased to accept them.

\begin{principle}[No new information]\label{p:noinfo}
Performing a thought experiment provides the thinker with no information about the world that she did not possess before it.
\end{principle}

\begin{principle}[Transformation]\label{p:trans}
If a process takes information a thinker possesses and yields a justified belief in a conclusion, the conclusion's justification is the justification the information confers, transmitted by the process; and the process transmits justification to the degree that it is reliable as a transition from that information to that conclusion, which is the degree to which it is good as an inference.
\end{principle}

Principle~\ref{p:noinfo} distinguishes thought experiments from real ones. A real experiment yields information from the world; the imagination receives none during a thought experiment. Principle~\ref{p:trans} says that a process that only transforms possessed information is, epistemically, an inference, whatever its psychological realization.

\section{The argument}

\begin{proposition}\label{p:encode}
Every thought experiment that justifies its conclusion is encoded by an argument.
\end{proposition}

\begin{proof}
Let a thought experiment justify its conclusion $C$. By Principle~\ref{p:noinfo}, the information on which the thinker's judgement draws was possessed before the thought experiment: her beliefs about bodies, forces, logic and the like. By Principle~\ref{p:trans}, the justification of $C$ is the justification conferred by the propositions she possesses on which her judgement depends, transmitted by a transition that is good as an inference. Take those propositions as premises and that transition as the inferential steps; by Definition~\ref{d:enc}, the result is an argument that encodes the thought experiment.
\end{proof}

\begin{theorem}\label{t:main}
Thought experiments can teach us about nature, and the evidential force of a thought experiment is exactly the evidential force of the argument it encodes.
\end{theorem}

\begin{proof}
Thought experiments can teach us about nature because arguments can: an argument from justified premises about nature, by good inference, yields a justified conclusion about nature, which may be one the thinker did not previously believe, as in Galileo's case. For the second claim, the force of the thought experiment does not exceed that of the encoded argument, since by Principle~\ref{p:noinfo} there is no further source of justification, and by Principle~\ref{p:trans} the justification it has is that of the premises transmitted by the transition the argument represents. It does not fall short, since the thought experiment draws on the same premises by the same transition. So the two forces are equal.
\end{proof}

\begin{corollary}\label{c:fail}
A thought experiment fails exactly when the argument it encodes has an unjustified premise or a bad inference.
\end{corollary}

The corollary yields a test. Thought experiments can mislead, and opposing thought experiments sometimes support contradictory conclusions; on our account the disagreement is located and resolved by reconstructing the arguments and examining their premises, as \citet{Norton2004} argues.

\section{Galileo's falling bodies}

Galileo's thought experiment encodes a reductio. Its premises are: (1) the natural speed of fall of a body is an increasing function of its weight alone, which is Aristotle's thesis taken as the hypothesis for reductio; (2) when a faster body is joined to a slower one, the slower retards the faster and the faster speeds the slower, so that the composite falls at a speed between theirs; (3) two bodies tied together form a body whose weight is the sum of their weights. From (1) and (3), the composite falls faster than the heavy body; from (2), it falls more slowly. So (1) is false. If the speed of fall is a function of weight alone, it is then a constant function, and bodies of different weight fall at the same rate.

The reconstruction makes visible what the vivid scenario conceals. Premise (2) is an empirical generalization about mediative effects, drawn from experience of bodies in contact, and premise (3) is a claim about how weight composes. The conclusion that all bodies fall at the same rate needs the further premise that speed of fall depends on weight alone, which is false in a resisting medium: a feather and a coin do not fall together in air. Galileo's conclusion holds for fall in a vacuum, where that premise is true, and the reconstruction explains both the force of the thought experiment and its limits.

\section{Objections}

\begin{objection}[Platonism]
\citet{Brown1991,Brown2004} holds that in Galileo's case we pass from a false theory to a new one without new data and without a valid deduction from old premises, since the conclusion that all bodies fall alike is not entailed by the refutation of Aristotle. The best explanation is that we grasp the law directly.
\end{objection}

\begin{reply}
The conclusion is entailed once the premise that speed depends on weight alone is made explicit, as the reconstruction shows. That this premise was tacit does not show that the conclusion was grasped without it: the restriction of Galileo's conclusion to the vacuum, where the premise holds, shows that it was in play. The platonist's faculty for perceiving relations among universals is not needed to explain any case, and could not explain why some thought experiments mislead.
\end{reply}

\begin{objection}[Tacit knowledge]
\citet{Gendler1998} argues that Galileo's thought experiment works by bringing the reader to recognize that whether two bodies tied by a string form one body or two is irrelevant to how they fall. This recognition is not a premise accepted in advance; it is produced by the imagined scenario.
\end{objection}

\begin{reply}
The recognition draws on what the reader already knows about bodies, strings and weight: imagining the scenario prompts her to apply that knowledge, and the knowledge, not the imagining, justifies the recognition. By Definition~\ref{d:enc}, a premise may be tacit, and the function of the scenario is to make tacit premises available to inference. Gendler's point shows that thought experiments are psychologically effective ways of eliciting premises, which we accept; it does not show that their evidential force exceeds that of the elicited premises.
\end{reply}

\begin{objection}[Mental models]
\citet{Nersessian1993} and \citet{Miscevic1992} hold that thought experimenters reason by simulating a scenario in a mental model, not by manipulating propositions. Thought experiments are therefore not arguments.
\end{objection}

\begin{reply}
The thesis is about evidential force, not psychological mechanism. A mental simulation that takes possessed information about physical behavior and yields a conclusion is, by Principle~\ref{p:trans}, epistemically an inference, and its reliability is the reliability of the transition it effects. Model-based reasoning is a way of performing inferences, as \citet{Norton2004} observes.
\end{reply}

\begin{objection}[Reconceptualization]
\citet{Kuhn1977} argued that some thought experiments teach not by deriving conclusions but by exposing incoherence in a scientist's concepts, forcing a conceptual change that no argument from her old premises could deliver.
\end{objection}

\begin{reply}
Exposing an incoherence is deriving a contradiction from the premises one accepts, and that is an argument, the reductio form of Galileo's case. The conceptual change that follows is justified, where it is, by further premises about which concepts best fit the phenomena; what the thought experiment itself establishes is that the old concepts are inconsistent.
\end{reply}

\section{Conclusion}

Thought experiments can teach us about nature, because they encode arguments from premises we already possess, and arguments can yield new conclusions about nature. Since a thought experiment supplies no new information, its evidential force is that of the encoded argument, neither more nor less. Galileo's falling bodies is a reductio whose force depends on tacit premises about how speed depends on weight; when those premises are made explicit, the thought experiment's force and its limits are both explained.

\begin{thebibliography}{9}\small
\bibitem[Brown(1991)]{Brown1991} Brown, J.~R. (1991). \emph{The Laboratory of the Mind: Thought Experiments in the Natural Sciences}. Routledge.
\bibitem[Brown(2004)]{Brown2004} Brown, J.~R. (2004). Why thought experiments transcend empiricism. In C.~Hitchcock (ed.), \emph{Contemporary Debates in Philosophy of Science}. Blackwell.
\bibitem[Gendler(1998)]{Gendler1998} Gendler, T.~S. (1998). Galileo and the indispensability of scientific thought experiment. \emph{The British Journal for the Philosophy of Science}, 49(3), 397--424.
\bibitem[Kuhn(1977)]{Kuhn1977} Kuhn, T.~S. (1977). A function for thought experiments. In \emph{The Essential Tension}. University of Chicago Press.
\bibitem[Mi\v{s}\v{c}evi\'c(1992)]{Miscevic1992} Mi\v{s}\v{c}evi\'c, N. (1992). Mental models and thought experiments. \emph{International Studies in the Philosophy of Science}, 6(3), 215--226.
\bibitem[Nersessian(1993)]{Nersessian1993} Nersessian, N.~J. (1993). In the theoretician's laboratory: Thought experimenting as mental modeling. In \emph{PSA 1992}, vol.~2. Philosophy of Science Association.
\bibitem[Norton(1996)]{Norton1996} Norton, J.~D. (1996). Are thought experiments just what you thought? \emph{Canadian Journal of Philosophy}, 26(3), 333--366.
\bibitem[Norton(2004)]{Norton2004} Norton, J.~D. (2004). On thought experiments: Is there more to the argument? \emph{Philosophy of Science}, 71(5), 1139--1151.
\end{thebibliography}

\end{document}
